\begin{table}[h!]
	\centering
	\small
	\caption{جمع‌بندی روش‌های مختلف انطباق در زمان آزمون برخط}
	\label{table:tta-comparison}
	\begin{tabular}{p{2.5cm} p{6cm} p{6cm}}
		\toprule
		\textbf{روش} & \textbf{مزایا} & \textbf{معایب} \\ 
		\midrule
		\lr{TENT}\cite{ftent} & پیاده‌سازی ساده؛ بدون نیاز به پارامترهای اضافی. & مستعد راه‌حل‌های بدیهی؛ استفاده مستقیم از آنتروپی و بهینه‌ساز آزاد بدون محدودیت در به‌روزرسانی وزن‌های مدل، می‌تواند منجر به فروپاشی مدل به‌ویژه در حالتهای پیوسته شود. \\ 
		\midrule
		\lr{SAR}\cite{sar} & تحلیل دقیق موارد شکست با توجه به تغییرات ناگهانی در اندازه‌ی گرادیان‌ها و طراحی تابع هدف بر اساس این مشاهدات که به طور موثر از فروپاشی مدل در سنارو‌های مختلف جلوگیری می‌کند.& نیازمند تنظیم دقیق لایه‌هایی که باید منجمد یا قابل آموزش باشند، زیرا لایه‌های پایانی معمولاً به تغییرات حساس‌اند؛ تعداد فرا‌پارامتر‌های روش زیاد است و لزوم انتخاب عمق، کران و شعاع اغتشاش مناسب، استفاده از این روش را در محیط‌های عملیاتی و برخط دچار چالش جدی می‌کند. \\ 
		\midrule
		\lr{DeYo}\cite{deyo} & تنها نمونه‌هایی که آنتروپی پایین و اختلاف معنی‌دار دارند وارد انطباق می‌شوند؛ تابع هزینه تمرکز روی نمونه‌های معنادار را تضمین می‌کند و از فروپاشی مدل جلوگیری می‌کند. & سربار محاسباتی ناشی از اعمال آشفتگی؛ نیاز به تنظیم دقیق آستانه‌ها؛ اگر مدل بر اساس ویژگی‌های محلی یک نمونه پیش‌بینی دقیقی ارائه دهد، اعمال آشفتگی پیش‌بینی مدل را تغییر نمی‌دهد و در این موارد این نمونه‌های مطمئن از فرایند انطباق حذف می‌شوند. \\ 
		\bottomrule
	\end{tabular}
\end{table}

\begin{table}[h!]\ContinuedFloat
	\centering
	\small
	\caption{جمع‌بندی روش‌های مختلف انطباق در زمان آزمون برخط}
	\begin{tabular}{p{2.5cm} p{6cm} p{6cm}}
		\toprule
		\textbf{روش} & \textbf{مزایا} & \textbf{معایب} \\ 
		\midrule
		\lr{TEA}\cite{tea} & استفاده از انرژی به جای آنتروپی؛ ترجیح عدم قطعیت به جای بیش‌اعتمادی؛ موثر زمانی که مدل ذاتاً نامطمئن است؛ کاهش خطای \lr{ECE} و \lr{MCE}. & کنترل مستقیمی روی تابع انرژی وجود ندارد؛  مشابه آنتروپی استفاده از تابع انرژي بدون اعمال ساز‌وکار‌های مناسب تشخیص و حذف نمونه‌های نامطمئن مدل را مستعد فروپاشی می‌کند؛ استفاده از \lr{SGLD} باعث کندی زیاد شده و استفاده از این روش را در محیط‌های بلادرنگ عملا ناممکن می‌کند. \\ 
		\midrule
		\lr{AEA}\cite{aea} & ترکیب انرژی و آنتروپی باعث تسریع انطباق و تعادل میان بیش‌اعتمادی و عدم قطعیت می‌شود؛ عملکرد بهتر در برخی مجموعه‌داده‌ها و مدل‌ها مشاهده شده است. & حساسیت بالا نسبت به هم‌ترازی  توابع انرژی و آنتروپی در تابع هزینه؛ نیاز به تنظیم دقیق پارامترهای $\lambda_1$ و $\lambda_2$ که در شرایط برخط بسیار چالش‌برانگیز است. \\ 
		\midrule
		\lr{T3A}\cite{t3a} & استفاده از مجموعه پشتیبانی و فیلتر آنتروپی باعث پایدار ماندن پیش‌بینی‌ها و جلوگیری از تاثیر نمونه‌های با برچسب موقت نادرست می‌شود. & وابستگی به انتخاب درست آستانه آنتروپی \(\alpha_k\)؛ احتمال اختصاص اشتباه برچسب موقت که عملکرد مدل را کاهش می‌دهد. \\ 
		\midrule
		\lr{SHOT}\cite{shot} & ترکیب مزایای روش‌های مبتنی بر آنتروپی با سازوکار پایدار برچسب‌زنی خودکار؛ انطباق سریع و کاهش اثر برچسب‌های دارای اغتشاش. & برای مسائل برخط و جریان داده‌ای طراحی نشده است و به کار گیری آن برای این دسته مسائل نیاز به ترکیب کردن آن با روش‌های دیگر دارد؛ راهکار مشخصی برای جلوگیری از فروپاشی مدل ارائه نمی‌کند؛ حساس به عدم توازن برچسب‌ها در نمونه‌ها. \\ 
		\bottomrule
	\end{tabular}
\end{table}

\begin{table}[h!]\ContinuedFloat
	\centering
	\small
	\caption{جمع‌بندی روش‌های مختلف انطباق در زمان آزمون برخط}
	\begin{tabular}{p{2.5cm} p{6cm} p{6cm}}
		\toprule
		\textbf{روش} & \textbf{مزایا} & \textbf{معایب} \\ 
		\midrule
		\lr{REM}\cite{rem} & حفظ استحکام و انعطاف‌پذیری مدل در زمان آزمون؛ کاهش حساسیت به تغییرات ناگهانی داده‌ها. & طول زنجیره و نسبت پوشش‌ها به عنوان فراپارامترهای کلیدی نیازمند تنظیم دقیق هستند؛ سازوکار مشخصی برای حذف نمونه‌های نامطمئن ارائه نمی‌دهد که در حالتهای برخط چالش‌برانگیز است. \\ 
		\midrule
		\lr{PETAL}\cite{petal} & کاهش اثر فراموشی فاجعه‌بار و انباشت خطا؛ انطباق پایدار روی داده‌های بدون برچسب؛ عدم نیاز به دسترسی به داده‌های منبع و حفظ دانش اولیه مدل. & پیچیدگی محاسباتی بالا؛ حساسیت به تنظیم فراپارامترها؛ محدودیت در جریان داده‌های بسیار تغییرپذیر؛ تقریب گوسی ممکن است تمامی پیچیدگی فضای وزن‌ها را مدل‌سازی نکند. \\ 
		\midrule
		\lr{COTTA}\cite{cotta} & جلوگیری از فروپاشی مدل با استفاده از معماری معلم-دانش‌آموز و میانگین‌گیری متحرک؛ ساده و مؤثر در کاهش انباشت خطا. & از نظر محاسباتی بسیار پرهزینه؛ انطباق کند؛ کاهش انعطاف ‌پذیری. \\ 
		\midrule
		\lr{EATA}\cite{eata} & ارائه روشی برای حذف نمونه‌های تکراری که سرعت انطباق را افزایش می‌دهد و به کاهش انباشت خطا و فراموشی فاجعه‌بار کمک می‌کند؛ تنها نیازمند چند گام پس‌انتشار روی نمونه‌های منتخب بدون سربار محاسباتی زیاد. & تکیه بر آنتروپی برای انتخاب نمونه‌های مطمئن که لزوماً معیار قابل اعتمادی برای حذف نمونه‌های دارای اغتشاش نیست و می‌تواند در مواردی فروپاشی مدل را تسریع کند. \\ 
		\bottomrule
	\end{tabular}
\end{table}